\chapter{The Floating-Point Model}
\label{app:fp}

This appendix collects background on floating-point arithmetic that is used throughout the course. Supplementary reading:
\begin{itemize}
  \item \cite{TB} Lectures 13 and 14.
  \item \cite{Demmel} Section 1.5.
  \item \cite{GVL} Section 2.4.
\end{itemize}

\section{IEEE 754 Formats}
\label{sec:ieee}

The formats in common use are listed below; $b_p$ counts the implicit bit, so a significand field with $b_p - 1$ stored bits provides $b_p$ bits of precision.
\begin{center}
\begin{tabular}{lcccccc}
  \toprule
  Format & Total bits & Exponent bits & Stored significand bits & $b_p$ & Bias & Decimal digits \\
  \midrule
  Float64 & 64 & 11 & 52 & 53 & 1023 & $\approx 15.9$ \\
  Float32 & 32 & 8 & 23 & 24 & 127 & $\approx 7.2$ \\
  Float16 & 16 & 5 & 10 & 11 & 15 & $\approx 3.3$ \\
  BFloat16 & 16 & 8 & 7 & 8 & 127 & $\approx 2.4$ \\
  \bottomrule
\end{tabular}
\end{center}
The exponent field uses a biased encoding, and its two extreme values are reserved:
\begin{itemize}
  \item An all-zeros exponent encodes $\pm 0$ (if the significand is also zero) or a \emph{subnormal} number, which drops the implicit leading $1$ and uses $0.f \times 2^{E_{\min}}$. Subnormals provide \emph{gradual underflow}, which guarantees $x \ne y \implies x - y \ne 0$.
  \item An all-ones exponent encodes $\pm\infty$ (zero significand) or \texttt{NaN} (nonzero significand).
\end{itemize}
The normal range of Float64 is therefore roughly $2.2 \times 10^{-308}$ to $1.8 \times 10^{308}$. IEEE defines five exceptions (invalid, division by zero, overflow, underflow, inexact), and specifies that $\infty$ and \texttt{NaN} propagate through arithmetic instead of halting it. Demmel~\cite[\S1.5]{Demmel} stresses the value of this: an algorithm can run to completion and check its result at the end, instead of branching at every step. Note that \texttt{NaN != NaN} is always true, so one must test with \texttt{isnan}.

\section{Machine Epsilon and the Model of Arithmetic}

\subsection{Two Conventions for Machine Epsilon}

Two different quantities are both called ``machine epsilon'' in the literature, and they differ by a factor of $2$.

\begin{definition}[Machine Epsilon and Unit Roundoff]
  \label{def:eps}
  Let $b_p$ be the precision of a floating-point format.
  \begin{itemize}
    \item The \emph{machine epsilon} $\eps = \epsm \doteq 2^{-(b_p - 1)}$ is the distance from $1$ to the next larger floating-point number. For Float64, $\eps = 2^{-52} \approx 2.22 \times 10^{-16}$. This is the convention of the lecture, of Julia's \texttt{eps(Float64)}, and of C's \texttt{DBL\_EPSILON}.
    \item The \emph{unit roundoff} $u \doteq 2^{-b_p} = \eps/2$ is the largest relative rounding error, i.e., half a spacing. For Float64, $u = 2^{-53} \approx 1.11 \times 10^{-16}$. Trefethen and Bau~\cite{TB} call this quantity $\epsilon_{\text{machine}}$, and Golub and Van Loan~\cite{GVL} call it the unit roundoff.
  \end{itemize}
\end{definition}

Any statement of the form $O(\eps)$ is unaffected by the choice, but constants copied from different textbooks must be read with care.

The IEEE default rounding mode is \emph{round to nearest, ties to even}, under which both orders in \Cref{ex:nonassoc} (Lecture 1) give $4$. Non-associativity does not depend on the rounding mode, however. With ties to even,
\begin{equation}
  (1 + u) + u = 1 \qquad \text{but} \qquad 1 + (u + u) = 1 + \eps,
\end{equation}
since on the left each intermediate result is a tie that rounds back to $1$, while on the right $u + u = \eps$ and $1 + \eps$ are both exact.

\subsection{The Fundamental Axiom}

Error analysis needs rounding to be captured by a rule we can reason with. Let $\fl(x)$ denote the floating-point number nearest to $x$. Because the relative spacing is bounded, every $x$ in the normal range satisfies
\begin{equation}
  \fl(x) = x(1 + \delta), \qquad \abs{\delta} \le u. \label{eq:fl}
\end{equation}
IEEE requires $+$, $-$, $\times$, $\div$, and $\sqrt{\cdot}$ to be \emph{correctly rounded}: the exact result is computed in infinite precision and then rounded once. This gives the following model, stated in the same form in \cite[Lecture 13]{TB}, \cite[\S2.4]{GVL}, and \cite[\S1.5]{Demmel}.

\begin{theorem}[Fundamental Axiom of Floating-Point Arithmetic]
  \label{thm:axiom}
  For all floating-point numbers $x, y$ and $\circledast \in \{+, -, \times, \div\}$,
  \begin{equation}
    \fl(x \circledast y) = (x \circledast y)(1 + \delta), \qquad \abs{\delta} \le u.
  \end{equation}
\end{theorem}

The axiom has an immediate \emph{backward} reading. For addition,
\begin{equation}
  \fl(x + y) = (x + y)(1 + \delta) = x(1 + \delta) + y(1 + \delta),
\end{equation}
so the computed sum is \emph{exactly} the sum of two slightly perturbed inputs. Pushing errors back onto the inputs in this way is the main thread of error analysis in this course; see \Cref{sec:backward}.

\subsection{Accumulation of Errors}

A single operation incurs an error of size $O(u)$. For $n$ operations, Golub and Van Loan write
\begin{equation}
  \gamma_n \doteq \frac{nu}{1 - nu} \approx nu \qquad (nu \ll 1).
\end{equation}

\begin{proposition}[Error in Inner Products]
  \label{prop:dot}
  For $\vec{x}, \vec{y} \in \R^n$,
  \begin{equation}
    \abs{\fl(\vec{x}^\T \vec{y}) - \vec{x}^\T \vec{y}} \le \gamma_n \sum_{i=1}^{n} \abs{x_i}\abs{y_i} = \gamma_n \abs{\vec{x}}^\T \abs{\vec{y}}.
  \end{equation}
  Similarly, for matrices, $\abs{\fl(AB) - AB} \le \gamma_n \abs{A}\abs{B}$ entrywise.
\end{proposition}

The \emph{shape} of this bound matters much more than its constant. The right-hand side involves $\abs{\vec{x}}^\T\abs{\vec{y}}$, not $\abs{\vec{x}^\T\vec{y}}$. If all terms have the same sign, the two agree and the relative error is $O(nu)$, which is fine. If instead $\vec{x}^\T\vec{y}$ is a small number obtained by cancelling many terms of both signs, then $\abs{\vec{x}}^\T\abs{\vec{y}} \gg \abs{\vec{x}^\T\vec{y}}$ and the relative error can be arbitrarily bad. The linear factor $n$ is almost never the real issue (it is a very pessimistic worst case, and random errors grow more like $\sqrt{n}$). The real issue is always \emph{cancellation}.

\section{Backward Stability}
\label{sec:backward}

The textbooks~\cite[Lecture 14]{TB} mostly use a stronger notion than the stability defined in \Cref{sec:stability} of Lecture 1.

\begin{definition}[Backward Stability]
  \label{def:backward-stable}
  An algorithm $\hat{f}$ for $f$ is \emph{backward stable} if for every $x$ there exists $\tilde{x}$ with
  \begin{equation}
    \hat{f}(x) = f(\tilde{x}) \qquad \text{and} \qquad \frac{\abs{\tilde{x} - x}}{\abs{x}} = O(u).
  \end{equation}
\end{definition}

In words: \emph{the computed answer is the exact answer to a nearby problem}. The backward reading of \Cref{thm:axiom} is the simplest example. Backward stability is the gold standard of numerical linear algebra because it cleanly separates the problem from the algorithm, and it yields the classical rule of thumb
\begin{equation}
  \text{forward relative error} \lesssim \kappa_f(x) \cdot u,
\end{equation}
or, in a form that is easier to remember,
\begin{equation}
  \text{accuracy} \approx \text{condition number (problem)} \times \text{backward error (algorithm)}.
\end{equation}
A backward stable algorithm still returns inaccurate answers for ill-conditioned problems, and this is \emph{not the algorithm's fault}; it cannot be improved upon. Conversely, when a result is inaccurate, the first step should be to estimate the condition number, not to suspect the code.

\begin{center}
\begin{tabular}{l p{0.2\textwidth} p{0.34\textwidth} l}
  \toprule
  Concept & Belongs to & Question it answers & Measure \\
  \midrule
  Conditioning & problem $f$ and input $x$ & how much input perturbations are intrinsically amplified & $\kappa_f(x)$ \\
  Stability & algorithm $\hat{f}$ & whether the implementation amplifies errors further & backward error $O(u)$ \\
  Accuracy & final result & whether the computed number is right & $\lesssim \kappa \cdot u$ \\
  \bottomrule
\end{tabular}
\end{center}

Some common condition numbers:
\begin{center}
\begin{tabular}{lll}
  \toprule
  Problem & $\kappa$ & Remark \\
  \midrule
  $f(x) = \sqrt{x}$ & $1/2$ & always well-conditioned; even shrinks errors \\
  $f(x) = x^n$ & $n$ & grows linearly with the power \\
  $f(x) = e^x$ & $\abs{x}$ & ill-conditioned for large $\abs{x}$ \\
  $f(x_1, x_2) = x_1 - x_2$ & $\dfrac{\abs{x_1} + \abs{x_2}}{\abs{x_1 - x_2}}$ & cancellation (componentwise perturbations) \\
  $\vec{x} \mapsto A\vec{x}$ & $\dfrac{\norm{A}\norm{\vec{x}}}{\norm{A\vec{x}}} \le \kappa(A)$ & depends on the direction of $\vec{x}$ \\
  solve $A\vec{x} = \vec{b}$ & $\kappa(A) = \norm{A}\norm{A^{-1}}$ & \\
  polynomial roots & can be enormous & Wilkinson's polynomial \\
  \bottomrule
\end{tabular}
\end{center}
For $f : \R^n \to \R^m$ one replaces $f'$ by the Jacobian $J(\vec{x})$: the absolute condition number is $\hat{\kappa} = \norm{J(\vec{x})}$ and the relative condition number is $\kappa = \norm{J(\vec{x})}\norm{\vec{x}}/\norm{f(\vec{x})}$~\cite[Lecture 12]{TB}.

